\begin{theorem}[Reflection to the base quasi-order]
Let $Q$ be a quasi-order.  If
$h:\mathsf{IncSeq}\to\mathsf{PowerQ}(Q)$ is locally constant and bad for
$\mathsf{PowerRel}(r)$, then there is a locally constant multi-sequence
$g:\mathsf{IncSeq}\to Q$ which is bad for $r$ and satisfies
\[
g(x)\in\mathsf{PowerSupport}(h(x))
\qquad\text{for every }x\in\mathsf{IncSeq}.
\]
\end{theorem}
\begin{proof}
For an increasing enumeration $z$, write
\[
L_{n,j}(z)=\mathsf{PowerCopiedLeft}_r(h,z,n,j).
\]
The copied-row definition \noderef{PowerCopiedLeft} says that at depth zero
$L_{n,0}(z)$ is the canonical response \noderef{PowerIResponse} to the
adjacent pair
\[
\bigl(h(\mathsf{PowerShiftIter}(n,z)),
      h(\mathsf{PowerShiftIter}(n+1,z))\bigr),
\]
and that at depth $j+1$ the value in row $n$ is the same canonical response
to $L_{n,j}(z)$ and $L_{n+1,j}(z)$.  Thus row $n+1$ supplies the copied right
move, and all later values are determined by these fixed responses.

Fix $x$.  Applying the termination statement
\noderef{PowerCopiedTerminates} with $n=0$ shows that this copied play
reaches a depth satisfying the terminal condition
\noderef{PowerCopiedTerminalAt}, so the first two row values are atoms.  We choose the
common depth and labels supplied by the finite-terminal conclusion
\noderef{PowerCopiedFiniteTerminal}; write them as $k,q_x,q'_x$, so that
\[
L_{0,k}(x)=\mathsf{atom}(q_x),\qquad
L_{1,k}(x)=\mathsf{atom}(q'_x).
\]
The same conclusion says that if a replacement multi-sequence agrees with
$h$ on the shift values through index $k+2$, then its first-row value at
depth $k$ is still $\mathsf{atom}(q_x)$.

The first conjunct of \noderef{PowerCopiedFailureMove}, at row $0$ and
depth $k$, gives
\[
\neg\mathsf{PowerRel}(r)(L_{0,k}(x),L_{1,k}(x)).
\]
Substituting the two atom equations and using the atom--atom equivalence
\noderef{PowerRelAtomAtom} yields
\[
\neg r(q_x,q'_x). \tag{1}
\]
The shift identity \noderef{PowerCopiedShiftIdentity}, with $n=0$, gives
\[
L_{0,k}(\mathsf{ShiftMap}(x))=L_{1,k}(x)
  =\mathsf{atom}(q'_x). \tag{2}
\]

We next make the terminal label independent of all choices.  The second
conjunct of \noderef{PowerCopiedFailureMove} says that each successive
$L_{n,j+1}(z)$ is a legal move from $L_{n,j}(z)$.  By the atom clause in the
definition \noderef{PowerMove}, a legal move whose source is
$\mathsf{atom}(q)$ must equal $\mathsf{atom}(q)$.  Hence, once a row value is
an atom, it remains that same atom at every later depth.  If two depths in
the same row are terminal, compare them in increasing order and use this
persistence; the two atom labels are equal.  By
\noderef{PowerCopiedTerminates}, every row-zero copied play has a terminal
depth, so we may define $g(z)$ to be its unique terminal row-zero atom.
Equation (2) and the same uniqueness give
\[
g(x)=q_x,\qquad g(\mathsf{ShiftMap}(x))=q'_x. \tag{3}
\]

For the support assertion, \noderef{PowerCopiedSupport} at $n=0$ and
depth $k$ gives
\[
\mathsf{PowerSupport}(L_{0,k}(x))
 \subseteq \mathsf{PowerSupport}(h(\mathsf{PowerShiftIter}(0,x))).
\]
By the atom clause of \noderef{PowerSupport}, the left side contains $q_x$;
by the defining zeroth-iterate equation in \noderef{PowerShiftIter}, the
right side is $\mathsf{PowerSupport}(h(x))$.  Therefore
\[
g(x)=q_x\in\mathsf{PowerSupport}(h(x)). \tag{4}
\]
Combining (1) and (3), for every $x$ we have
\[
\neg r\bigl(g(x),g(\mathsf{ShiftMap}(x))\bigr).
\]
By the definition \noderef{BadMultiSequence}, this is exactly
$\mathsf{BadMultiSequence}\,r\,g$.

It remains to prove local constancy.  Fix $x$, and use the finite-terminal
choice above, with its depth $k$ and label $q_x$.  Apply
\noderef{PowerFiniteOrbitConstancy} to $h$, $x$, and $m=k+2$.  We obtain a
finite prefix length $N$ such that, whenever
$\mathsf{PrefixAgree}(N,x,y)$, then
\[
h(\mathsf{PowerShiftIter}(j,y))
 =h(\mathsf{PowerShiftIter}(j,x))
\quad\text{for every }j\leq k+2. \tag{5}
\]
For such a $y$, (5) is exactly the hypothesis of
\noderef{PowerCopiedOrbitDependence} with row $n=0$ and depth $k$; the bound
there is $0+k+2$, with the equality orientation reversed by symmetry if
necessary.  Consequently,
\[
L_{0,k}(y)=L_{0,k}(x)=\mathsf{atom}(q_x). \tag{6}
\]
Apply \noderef{PowerCopiedTerminates} to $y$ and row $0$, and use the
terminal condition \noderef{PowerCopiedTerminalAt} to choose a depth $\ell$
whose first-row value is $\mathsf{atom}(p)$.  If $\ell\leq k$,
the legal moves supplied by \noderef{PowerCopiedFailureMove} carry this
atom unchanged at each successor depth, and induction over the finite
interval from $\ell$ to $k$ gives $p=q_x$.  If
$k\leq\ell$, induction over the finite interval carries the atom in (6) to
depth $\ell$ unchanged and again gives $p=q_x$.  Thus the unique terminal
label of $y$ is $q_x$, and
\[
g(y)=q_x=g(x).
\]
This holds for every $y$ with the finite-prefix agreement in (5).  By the
definition \noderef{LocallyConstant}, $g$ is locally constant, completing
the proof.
\end{proof}
